\section*{Capítulo \ref{ch:b3:green-industrial} --- Química verde e processos industriais}

\begin{solution}{exo:b3:green-industrial:1}
Diels--Alder: 100\,\% (uma adição). Esterificação: $88.0/(60.0 + 46.0) = 83$\,\%. Wittig:
$104.0/(106.0 + 276.0) = 27$\,\%, com o óxido de trifenilfosfina levando embora a maior parte da massa.
\end{solution}

\begin{solution}{exo:b3:green-industrial:2}
$\mathrm{PMI} = 120/8.0 = 15$; $E = \mathrm{PMI} - 1 = 14$.
\end{solution}

\begin{solution}{exo:b3:green-industrial:3}
Uma lavagem de uma carga padrão a uma dada temperatura e com um dado resultado de limpeza. Os detergentes
são dosados de modo diferente: um pó concentrado exige menos por lavagem, de modo que massas iguais
não prestam o mesmo serviço.
\end{solution}

\begin{solution}{exo:b3:green-industrial:4}
2-Metiltetra-hidrofurano ou acetato de etila no lugar do diclorometano; heptano (com acetato
de etila) no lugar do hexano; no lugar da DMF, acetato de etila, 2-metiltetra-hidrofurano ou
carbonato de dimetila, ou um acoplamento em água com \omterm{def:b3:colloids:surfactant}{tensoativos}, conforme os
reagentes.
\end{solution}

\begin{solution}{exo:b3:green-industrial:5}
1.000 mol de ácido; 0.800 mol de éster (\qty{88.0}{g/mol}): rendimento 80\,\%. AE 83.0\,\%. SF $= 152/106 = 1.43$.
$\mathrm{RME} = 70.4/152 = 0.463$; $\mathrm{AE} \times \text{rendimento}/\mathrm{SF} = 0.830 \times 0.80/1.43 = 0.463$.
\end{solution}

\begin{solution}{exo:b3:green-industrial:6}
Cloridrina: $44.0/(28.0 + 71.0 + 74.1) = 25$\,\%; oxidação direta: 100\,\%. Um mol de $\ce{CaCl2}$ por mol de
óxido de etileno: $\qty{1.0e6}{g}/\qty{44.0}{g/mol} \times \qty{111.1}{g/mol} = \qty{2.5}{t}$ por tonelada.
\end{solution}

\begin{solution}{exo:b3:green-industrial:7}
$3/8 \times \qty{5.88e4}{mol} \times \qty{44.0}{g/mol} = \qty{0.97}{t}$ de $\ce{CO2}$ por tonelada de amônia. Para o hidrogênio, $\ce{CH4 + 2H2O -> CO2 + 4H2}$:
$1/4$ mol por mol, $\qty{5.0e5}{mol}/4 \times \qty{44.0}{g/mol} = \qty{5.5}{t}$ de $\ce{CO2}$ por tonelada de hidrogênio.
\end{solution}

\begin{solution}{exo:b3:green-industrial:8}
% ledger: dfg:H2O_l, dfh:H2O-l
Um quilograma corresponde a \qty{500}{mol}. A partir de $\Delta_rG^\circ = \qty{237.1}{kJ/mol}$: $\qty{1.19e5}{kJ} = \qty{32.9}{kWh}$. A partir de $\Delta_rH^\circ = \qty{285.8}{kJ/mol}$: \qty{39.7}{kWh}
(a diferença é calor que o meio externo pode fornecer).
\end{solution}

\begin{solution}{exo:b3:green-industrial:9}
$\qty{1.0e6}{g}/\qty{146.0}{g/mol} = \qty{6.85e3}{mol}$ de $\ce{N2O}$, $\qty{0.30}{t}$; $\times 273 = \qty{82}{t}$ de $\ce{CO2}$ equivalente por tonelada de ácido adípico.
\end{solution}

\begin{solution}{exo:b3:green-industrial:10}
A RME de uma etapa conta o reagente em excesso como consumido; se o excesso é recuperado
e reintroduzido, só a parte realmente perdida é consumida, e a eficiência mássica
global é maior que o produto dos valores das etapas. No \omterm{def:b3:green-industrial:e-factor}{fator E}, a
corrente reciclada não é resíduo, desde que a reciclagem esteja dentro da fronteira;
sua energia e suas perdas precisam então ser contadas.
\end{solution}

\begin{solution}{exo:b3:green-industrial:11}
Antes: $20 \times 0.80 = \qty{16}{kg}$ de solvente, 10\,\% perdidos: \qty{1.6}{kg} de resíduo. Depois: \qty{4.0}{kg}, \qty{0.40}{kg} perdidos. O
\omterm{def:b3:green-industrial:e-factor}{fator E} completo cai 1.2 kg por quilograma de produto (e a energia de
destilação, três quartos).
\end{solution}

\begin{solution}{exo:b3:green-industrial:12}
Por \omterm{def:b3:green-industrial:lca}{unidade funcional} e com a mesma fronteira, ela apresentaria cada categoria de impacto
separadamente (clima, uso da terra, uso de água\dots), em vez de um único escore. Uma
decisão exige os pesos atribuídos a essas categorias, a escassez local de água
e de terra, a incerteza dos dados e saber se a matéria-prima de base biológica
compete com a alimentação.
\end{solution}

\begin{solution}{pb:b3:green-industrial:1}
\textbf{1.} $\ce{C10H14 + C4H6O3 -> C12H16O + C2H4O2}$ (HF); $\ce{C12H16O + H2 -> C12H18O}$ (níquel de Raney);
$\ce{C12H18O + CO -> C13H18O2}$ (paládio).

\textbf{2.} $\ce{C10H14 + C4H6O3 + H2 + CO -> C13H18O2 + C2H4O2}$.

\textbf{3.} $134.0 + 102.0 + 2.0 + 28.0 = \qty{266.0}{g/mol}$ de reagentes; ibuprofeno \qty{206.0}{g/mol}: AE $= 206.0/266.0 = 77$\,\%.

\textbf{4.} Isobutilbenzeno, anidrido acético, cloroacetato de etila, etóxido de sódio,
ácido aquoso ($\ce{H3O+}$), hidroxilamina e água.

\textbf{5.} $134.0 + 102.0 + 122.5 + 68.0 + 19.0 + 33.0 + 2 \times 18.0 = \qty{514.5}{g/mol}$: AE $= 206.0/514.5 = 40$\,\%.

\textbf{6.} Ácido acético, cloreto de sódio, etanol, dióxido de carbono, água, amônia (na forma de
sal de amônio) e sais de alumínio provenientes do cloreto de alumínio estequiométrico.

\textbf{7.} $0.71 + 0.55 + 0.010 + 0.14 + 0.05 = \qty{1.46}{t}$.

\textbf{8.} PMI 1.46; $E = 0.46$.

\textbf{9.} $0.46 - 0.31 = 0.15$.

\textbf{10.} PMI $= 2.0 + 1.5 = 3.5$; $E = 2.5$.

\textbf{11.} Rendimentos abaixo de 100\,\%, excessos de reagentes, perdas de solvente e de catalisador e
materiais de isolamento acrescentam resíduos que a \omterm{def:b3:green-industrial:atom-economy}{economia atômica} ignora.

\textbf{12.} Prevenir os resíduos, maximizar a \omterm{def:b3:green-industrial:atom-economy}{economia atômica}, usar catalisadores em vez de
reagentes estequiométricos, evitar derivados e etapas desnecessários, economizar energia.

\textbf{13.} O cloreto de alumínio estequiométrico, hidrolisado em resíduo de alumínio no
isolamento; o fluoreto de hidrogênio é ao mesmo tempo catalisador e solvente, e é recuperado e reutilizado.

\textbf{14.} O níquel de Raney (um catalisador de hidrogenação heterogêneo).

\textbf{15.} Um complexo fosfina--paládio: a química de \omterm{def:b3:homogeneous-catalysis:carbonylation}{carbonilação} do
\cref{ch:b3:homogeneous-catalysis}.

\textbf{16.} O anidrido acético transfere um grupo acetila; a outra metade da
molécula sai como ácido acético.

\textbf{17.} O fluoreto de hidrogênio é muito tóxico e corrosivo, o monóxido de carbono, tóxico e
inflamável: equipamentos fechados de materiais resistentes, detectores e operadores
treinados.

\textbf{18.} Cada etapa acrescenta aquecimento, resfriamento, separações e recuperação de solvente; três
etapas exigem menos de tudo isso.

\textbf{19.} $2.5 \times 3000 = \qty{7500}{t}$ por ano.

\textbf{20.} $0.15 \times 3000 = \qty{450}{t}$ por ano.

\textbf{21.} Cerca de \qty{7000}{t} de resíduos evitados por ano.

\textbf{22.} Não: a natureza dos resíduos importa (uma tonelada de água salgada não é uma tonelada de
um composto tóxico persistente), assim como a energia, as emissões e os perigos dos
reagentes.

\textbf{23.} Os impactos da produção dos reagentes e da energia a montante, das emissões para
o ar e a água, e do uso e do descarte do produto, por \omterm{def:b3:green-industrial:lca}{unidade funcional}, em
várias categorias de impacto.

\textbf{24.} A rota catalítica em três etapas tem uma \omterm{def:b3:green-industrial:atom-economy}{economia atômica} de cerca de 77\,\% (40\,\% para a
rota em seis etapas).
\end{solution}
